\documentclass[10pt]{article}
\usepackage[margin=0.7in]{geometry}
\usepackage{amsmath}
\usepackage{array}
\usepackage{booktabs}
\usepackage{enumitem}
\usepackage[T1]{fontenc}
\usepackage{lmodern}
\setlength{\parindent}{0pt}
\setlength{\parskip}{3pt}
\begin{document}
\section*{Sprint 3 Update}
\subsection*{Sprint Goal}
Validate the rotor aerodynamic model and assess the numerical optimum
using active-constraint, KKT, and sensitivity analyses.
\subsection*{Planned and Completed Work}
\begin{table}[htbp]
\centering
\renewcommand{\arraystretch}{1.05}
\setlength{\tabcolsep}{4pt}
\begin{tabular}{@{}p{0.82\linewidth}p{0.10\linewidth}@{}}
\toprule
\textbf{Task} & \textbf{Owner} \\
\midrule
Benchmark BEMT predictions against published propeller measurements. & Ashish \\
Check active constraints, KKT regularity, and coupled-input sensitivity. & Ashish \\
\bottomrule
\end{tabular}
\end{table}
\subsection*{Technical Progress and Results}
\begin{itemize}[leftmargin=*,nosep]
\item Compared BEMT with published APC 10$\times$7 propeller
data. Static mean absolute percentage errors were 37.1\% for thrust
coefficient $C_T$ and 52.3\% for power coefficient $C_P$.
\item Examined the SLSQP solution using KKT and active-set concepts.
Binding limits were $c_{\min}$, $\theta_{75,\max}$, and $FM_{\min}$;
the active-constraint Jacobian had full rank.
\item Reoptimized nine provisional aircraft-mass/cruise-drag scenarios
(109.85--172.26 Wh). Refining the BEMT mesh from 96 to 192
elements changed the optimized screening energy by about 0.10\%.
\end{itemize}
\subsection*{Impediments and Lessons Learned}
Large experimental errors indicate that the assumed airfoil polar
is not yet suitable for physical rotor prediction. KKT checks
support local numerical consistency, not global optimality.
\subsection*{Next Sprint Plan}
Incorporate measured airfoil polars and rotor mass estimates;
update the coupled optimization using wing and propulsion inputs.
\end{document}
